\section{A boundary system for scalar curvature}
\label{four:sec:system}

Fix the prepared exterior and the geometric choices of
Proposition~\ref{four:prop:geometric-preparation}.  All contractions and
derivatives in this section use its background metric $g$, unless another
metric is indicated.  In particular, $\tau=\tr_gK$ is unrestricted.
The boundary normal $\nu$ points into $\Omega$.  Let $\Omega_R$ denote
the truncation at a large coordinate sphere $S_R$, so its ordinary outward
normal at its inner boundary $B$ is $-\nu$.

The purpose of the system is to produce a metric
$\ghat=e^{2t}(g+l(t)^2\dd f^2)$ with nonnegative scalar curvature and
negative boundary mean curvature.  The parameter $\eta_N$ below makes
the prescribed trace depend strictly increasingly on $f$.  A separate penalty prevents
$t$ from reaching $-\eps$.  This section proves that prevention using a
coarse height bound.  The narrower height interval needed to use
\eqref{four:eq:trace-slack} will be proved subsequently.

\subsection{Variables and a scalar-curvature identity}

The warped-graph scalar identity and coupled-Penrose strategy of Bray
and Khuri~\cite{BrayKhuri2010,BrayKhuri2011} extend the earlier
Jang--Schoen--Yau graph method~\cite{Jang1978,SchoenYau1981}.
The computation below fixes the dimension-four coefficients and
retains the unrestricted trace of $K$. The trace-discrepancy term
already appears before imposing the generalized Jang equation
in~\cite[Identity~9, p.~580]{BrayKhuri2011}. The prescribed $F=h+C$,
nonzero divergence source, and floor penalty are part of the
present system; its existence is proved independently, not inferred
from an older zero-divergence coupling.

Fix $0<\eps<1$ and a smooth nonincreasing function
$\vartheta:\R\to[0,1]$ equal to one on $(-\infty,0]$ and to zero on
$[1,\infty)$.  For every integer $N\ge4$ define
\begin{equation}\label{four:sys:parameters}
\begin{gathered}
 p(t)=N\vartheta(Nt),\qquad
 l(t)=\exp\left(\int_0^t p(s)\,\dd s\right),\\
 L_0(t)=e^{2t}l(t),\qquad
 \ell=e^{-\eps N},\qquad \eta_N=\ell^{3/2}.
\end{gathered}
\end{equation}
Thus $0\le p\le N$, $l(t)=e^{Nt}$ for $t\le0$, and $l$ is constant
with value between one and $e$ for $t\ge1/N$.  In particular
$l(-\eps)=\ell$.  The unknowns are real functions $f,Z$.  Define $t$
implicitly, and then all the remaining variables, by
\begin{equation}\label{four:sys:variables}
\begin{gathered}
 \sigma=|\grad f|,\qquad
 D=1+l(t)^2\sigma^2,\qquad d=D^{-1},\qquad
 Z=t+\frac16\log D,\qquad h=\eta_N f,\\
 w=\frac{l\grad f}{\sqrt D},\qquad a=|w|,\qquad
 v=a^2=1-d,\qquad \chi=\frac{3d}{3+pv},\\
 \Ad=\Id-w\otimes w,\qquad
 \Achi=\Id-\frac{3+p}{3+pv}w\otimes w,\qquad
 u=L_0\sqrt D.
\end{gathered}
\end{equation}
We identify vectors and covectors using $g$.  A symmetric endomorphism
also denotes the corresponding contravariant tensor, and
$|\xi|_A^2=A(\xi,\xi)$ for a covector $\xi$.
For fixed $\dd f$, the derivative of the defining expression for $Z$
with respect to $t$ is $1+pv/3\ge1$.  That expression tends to
$-\infty$ and $+\infty$ at the two ends of the real line.  It therefore
defines a unique smooth $t=t(x',Z,\dd f)$, with $x'$ the base point, for all inputs, including inputs
below the proposed floor.  No gradient bound is needed for this definition.
The eigenvalues of $\Ad$ and $\Achi$ on the line spanned by $\grad f$
are $d$ and $\chi$; their other three eigenvalues are one.  Consequently
\begin{equation}\label{four:sys:eigenvalues}
 0<\frac{3d}{3+N}\le\chi\le d\le1.
\end{equation}

Set
\begin{equation}\label{four:sys:metric-and-trace}
\begin{gathered}
 \gbar=g+l^2\dd f^2,\qquad \ghat=e^{2t}\gbar,\qquad
 H^f=\frac{l}{\sqrt D}\Hess_g f,\qquad S_f=K+H^f,\\
 F=\tr_{\Achi}S_f,\qquad
 V=u\Achi\bigl(3\grad Z+K(w,\cdot)\bigr).
\end{gathered}
\end{equation}
Where $\sigma>0$, let $e=\grad f/\sigma$.  Decompose $S_f$ into its
transverse block $T=S_f|_{e^\perp}$, mixed block
$M=S_f(e,\cdot)|_{e^\perp}$, and axial entry $j=S_f(e,e)$; write
$T^0=T-\frac13(\tr T)g|_{e^\perp}$.  Thus $\tr T=F-\chi j$.
Also write $x=e(t)$ and $y=(\dd t)|_{e^\perp}$.  The symbol $M$ in this
block notation is a covector, not the original manifold.

\begin{proposition}[Scalar identity]\label{four:sys:scalar-proposition}
For arbitrary smooth $f,Z$ and arbitrary symmetric $K$,
\begin{equation}\label{four:eq:scalar-identity}
 \frac12 e^{2t}\Scal_{\ghat}
 =\mu+J(w)+\cT+w(F)-F\tau-u^{-1}\divg_g V,
\end{equation}
where the smooth scalar $\cT$ is given, on $\{\dd f\ne0\}$, by
\begin{equation}\label{four:eq:quadratic-form}
\begin{aligned}
 \cT={}&\frac12|T^0|^2+|M+(p+1)ay|^2
       -v|y|^2+3(p+1)(|y|^2+d x^2)\\
 &-\frac13(F-\chi j)^2+\chi j^2-\chi Fj+F^2
       +2(p+1)\chi jax.
\end{aligned}
\end{equation}
The expression extends smoothly across $\{\dd f=0\}$, where any unit
axis may be used in the block expression.
\end{proposition}

\begin{proof}
We derive the identity in a stationary Lorentzian metric, keeping its
normal second form distinct from the prescribed tensor $K$.  On the
product with coordinate $s'$ use
\[
 \mathbf g=-l^2(\dd s'-\dd f)^2+\gbar.
\]
Its constant-$s'$ slices have metric $g$.  Their lapse, shift, and normal
second form, with the convention that the second form is half the normal
metric rate, are
\begin{equation}\label{four:sys:stationary-data}
 \begin{gathered}
 U=l\sqrt D,\qquad \beta=l^2\grad f=Uw,\\
 k=-\frac{\operatorname{sym}\nabla\beta}{U}
   =-H^f-p(\dd t\otimes w+w\otimes\dd t).
 \end{gathered}
\end{equation}
Here $\operatorname{sym}$ includes the factor $1/2$.  Put
$\theta=\tr_gk$, and let $N_0$ be the future unit normal.  The contracted
Gauss and normal expansion identities are
\[
 \Scal_g=\mathbf R+2\mathbf{Ric}(N_0,N_0)
                 +|k|^2-\theta^2,\qquad
 \mathbf{Ric}(N_0,N_0)=-N_0\theta-|k|^2+U^{-1}\Delta_gU.
\]
Stationarity gives $UN_0\theta=-\beta(\theta)$ and
$\divg\beta=-U\theta$.  If
\[
 \mathsf B(A,B)=A:B-(\tr_gA)(\tr_gB),\qquad
 P_A=A-(\tr_gA)g,
\]
the preceding two formulas consequently yield
\begin{equation}\label{four:sys:stationary-scalar}
 U\mathbf R=U\bigl(\Scal_g+\mathsf B(k,k)\bigr)
                   -2\divg(\grad U+\theta\beta).
\end{equation}
In the static coordinate $s'-f$ one instead has
\[
 \mathbf R=\Scal_{\gbar}-2l^{-1}\Delta_{\gbar}l,\qquad
 \Delta_{\gbar}\phi=\frac lU\divg\left(\frac Ul\Ad\grad\phi\right).
\]
Direct differentiation of $U=l\sqrt D$ gives
\[
 \grad U-\frac Ul\Ad\grad l=-k(\beta,\cdot).
\]
Moreover, using $J=\divg P_K$, symmetry of $P_K$, and
$\operatorname{sym}\nabla\beta=-Uk$ gives
\[
 \divg(P_K\beta)=U\{J(w)-\mathsf B(K,k)\}.
\]
Substitute these identities into \eqref{four:sys:stationary-scalar}, use
$\mu=(\Scal_g-\mathsf B(K,K))/2$, and put $Q=K-k$.  The result is
\begin{equation}\label{four:sys:bar-scalar}
 \frac12\Scal_{\gbar}
  =\mu+J(w)+\frac12\mathsf B(Q,Q)
       -U^{-1}\divg(UP_Qw).
\end{equation}
For example, the divergence arising directly from
\eqref{four:sys:stationary-scalar} is $U^{-1}\divg(UP_kw)$;
adding $J(w)$ combines it with $-U^{-1}\divg(UP_Kw)$ to give the
last term in \eqref{four:sys:bar-scalar}.  This also fixes its sign.

The four-dimensional conformal formula is
\[
 \frac12e^{2t}\Scal_{\ghat}
  =\frac12\Scal_{\gbar}-3\Delta_{\gbar}t-3|\dd t|_{\Ad}^2.
\]
Since
$\Delta_{\gbar}t=U^{-1}\divg(U\Ad\grad t)-p|\dd t|_{\Ad}^2$
and $u=e^{2t}U$, changing the divergence weight from $U$ to $u$ gives
\begin{equation}\label{four:sys:pretrace-identity}
\begin{aligned}
 \frac12e^{2t}\Scal_{\ghat}
 ={}&\mu+J(w)+\frac12\mathsf B(Q,Q)+2P_Q(w,\grad t)
               +3(p+1)|\dd t|_{\Ad}^2\\
 &-u^{-1}\divg\bigl(u(P_Qw+3\Ad\grad t)\bigr).
\end{aligned}
\end{equation}
The following differential identities follow from \eqref{four:sys:variables}:
\begin{equation}\label{four:sys:differential-identities}
\begin{aligned}
 3\dd Z&=(3+pv)\dd t+H^f(w,\cdot),\\
 \divg w&=\tr_{\Ad}H^f+p d\,w(t),\\
 \dd\log u&=(p+2+pv)\dd t+H^f(w,\cdot).
\end{aligned}
\end{equation}
They imply
\begin{equation}\label{four:sys:trace-divergence}
\begin{aligned}
 \frac Vu&=P_Qw+3\Ad\grad t+Fw,\\
 u^{-1}\divg(uw)
  &=\tr_gH^f+2(p+1)w(t)\\
  &=F+(1-\chi)j-\tau+2(p+1)w(t).
\end{aligned}
\end{equation}
Adding $uFw$ inside the divergence in
\eqref{four:sys:pretrace-identity} therefore leaves
$w(F)+F u^{-1}\divg(uw)$.  Its trace contribution is exactly
$-F\tau$, while the remaining quadratic expression is
\begin{equation}\label{four:sys:invariant-quadratic}
\begin{aligned}
 \cT={}&\frac12\mathsf B(Q,Q)+2P_Q(w,\grad t)
          +3(p+1)|\dd t|_{\Ad}^2\\
 &+F\left(F+\frac{3+p}{3+pv}S_f(w,w)+2(p+1)w(t)\right).
\end{aligned}
\end{equation}
This is smooth even at $w=0$, since
$(1-\chi)j=(3+p)S_f(w,w)/(3+pv)$.  For the block expansion, $Q$ has
blocks $T$, $M+pay$, $j+2pax$, and $\tr T=F-\chi j$.
Expanding \eqref{four:sys:invariant-quadratic} gives
\eqref{four:eq:quadratic-form}.  At $w=0$ it reduces to
\[
 \cT=\frac12\bigl(|S_f|^2+(\tr_gS_f)^2\bigr)
                      +3(p+1)|\dd t|^2,
\]
which is independent of the auxiliary axis.  This proves all assertions.
\end{proof}

\subsection{Coercivity and the boundary identity}

Throughout the proof, $\Pi_N$ denotes a bound of the form
$C(1+N)^{a_{\mathrm{pol}}}$ for a finite nonnegative integer $a_{\mathrm{pol}}$.  Its coefficient $C$ and
exponent $a_{\mathrm{pol}}$ may depend on the fixed prepared data, $\eps$, the chosen
collars and smooth switches, and any explicitly fixed small absorption
constant.  They never depend on $R$, a homotopy parameter, or a solution.
Different occurrences may denote larger such bounds.  Constants denoted
by $C(N)$ in later fixed-$N$ analytic estimates have no polynomial
restriction and will not enter a limit requiring polynomial control.

\begin{lemma}[Polynomial coercivity]\label{four:sys:coercivity-lemma}
For all inputs in \eqref{four:sys:variables},
\begin{equation}\label{four:eq:coercivity}
 |T|^2+|M|^2+d j^2+F^2+|\dd t|_{\Ad}^2\le\Pi_N\cT.
\end{equation}
At a point where $\dd t=0$, there are absolute constants $c,C>0$ such
that
\begin{equation}\label{four:sys:stationary-coercivity}
 \begin{aligned}
 c\bigl(|T^0|^2+|M|^2+\chi j^2+F^2\bigr)&\le\cT\\
 &\le C\bigl(|T^0|^2+|M|^2+\chi j^2+F^2\bigr).
 \end{aligned}
\end{equation}
\end{lemma}

\begin{proof}
Completing squares in \eqref{four:eq:quadratic-form} gives the exact identity
\begin{equation}\label{four:sys:exact-squares}
\begin{aligned}
 \cT={}&\frac12|T^0|^2+|M+(p+1)ay|^2
        +\bigl(3(p+1)-v\bigr)|y|^2\\
 &+3(p+1)d\left(x+\frac{\chi a j}{3d}\right)^2
       +\frac23\left(F-\frac{\chi j}{4}\right)^2
       +\frac{d(48-3d)}{8(3+pv)^2}j^2.
\end{aligned}
\end{equation}
For the final coefficient, the identity used in this completion is
\[
 \chi-\frac38\chi^2-\frac{(p+1)\chi^2v}{3d}
                 =\frac{d(48-3d)}{8(3+pv)^2}.
\]
The coefficient of $|y|^2$ is at least two, and the last coefficient is
at least $45d/[8(3+N)^2]$.  Thus $d j^2\le C(3+N)^2\cT$.
The inequalities $\chi^2\le d^2\le d$ then recover $F^2$ from its
shifted square and $d x^2$ from its shifted square, at polynomial cost.
Recovering $M$ costs at most $C(1+N)^2|y|^2$, and recovering $T$ uses
$\tr T=F-\chi j$.  This proves \eqref{four:eq:coercivity}; these observations,
for instance, allow one common bound $C(1+N)^2$ in that inequality.

When $\dd t=0$, complete only the $(F,j)$ terms to obtain
\begin{equation}\label{four:sys:stationary-squares}
 \cT=\frac12|T^0|^2+|M|^2
      +\frac23\left(F-\frac{\chi j}{4}\right)^2
      +\chi\left(1-\frac{3\chi}{8}\right)j^2.
\end{equation}
Since $5/8\le1-3\chi/8\le1$ and $\chi^2\le\chi$, this gives
\eqref{four:sys:stationary-coercivity} with absolute constants.
\end{proof}

\begin{lemma}[Boundary conversion]\label{four:sys:boundary-lemma}
On any smooth face on which $f$ is constant, let
$P_B=\tr_BK$, $w_\nu=\inner{w}{\nu}$, and let $H$ be its background
mean curvature for the normal into $\Omega$.  Then
\begin{equation}\label{four:eq:boundary-identity}
 \frac{V_\nu}{u}
   =d(H+3\partial_\nu t)-H+w_\nu(F-P_B),\qquad
 \widehat H=e^{-t}\sqrt d\,(H+3\partial_\nu t).
\end{equation}
\end{lemma}

\begin{proof}
With $\mathrm{II}_B(X,Y)=g(\nabla_X\nu,Y)$ for tangent vectors,
one has tangentially $\Hess f=(\partial_\nu f)\mathrm{II}_B$, so
$\tr T=P_B+w_\nu H$ and
$\chi j=F-P_B-w_\nu H$.  Use this relation in the first identity of
\eqref{four:sys:trace-divergence}; as $w$ is normal on the face, it gives
\[
 V_\nu/u=3d\partial_\nu t+w_\nu(F-P_B-w_\nu H),
\]
which is the first assertion because $w_\nu^2=1-d$.  The induced
metric of $\gbar$ on the face equals $g|_B$, and its normal in base
coordinates is $\sqrt d\,\nu$.  In the tangential first variation,
$l^2\dd f^2$ contributes zero since the tangential derivatives of $f$
vanish there.  Hence $H_{\gbar}=\sqrt d\,H$.  Applying the boundary
conformal formula with $\nu_{\gbar}=\sqrt d\,\nu$ proves the second
assertion.  Both assertions also hold when $\dd f=0$.
\end{proof}

\subsection{The equations and their four-stage homotopy}

Extend the end radius to a positive smooth function $r\ge1$ on $\Omega$
and set
\begin{equation}\label{four:sys:source}
\begin{gathered}
 \rho_0=r^{-4-\delta_1},\qquad \rho_1=r^{-3},\qquad
 m_0(t)=\vartheta\bigl(N(t+\eps)\bigr),\qquad
 \cP=C_N(\rho_0+v\rho_1),\\
 \Xi=\delta_0\cT+\rho+\delta_0\eta_N a\sigma-m_0(t)\cP.
\end{gathered}
\end{equation}
The constant $\delta_0>0$ and the polynomial $C_N$ will be fixed in
Proposition~\ref{four:prop:floor-exclusion}.  The smooth positive weight
$\rho$ comes from \eqref{four:eq:trace-slack}; in particular
$\rho\le C\rho_0$.  The combination
$a\sigma=l|\dd f|^2/\sqrt D$ is smooth also at $\dd f=0$, so all
terms of $\Xi$ are smooth functions of the indicated jets.
Choose a constant $N_B>1+\sup_B|H|$, and put
\[
 \cB=-H+w_\nu(F-P_B)-N_Bd.
\]
The desired system is
\begin{equation}\label{four:eq:system}
\begin{cases}
 F=h+C, &\text{in }\Omega_R,\\
 \divg V=u\Xi, &\text{in }\Omega_R,\\
 h=b,\quad V_\nu/u=\cB, &\text{on }B,\\
 f=Z=0, &\text{on }S_R.
\end{cases}
\end{equation}
In every boundary flux expression the symbol $F$ is evaluated using the
prescribed trace right side.  Thus the boundary condition contains no
second derivative of $f$.  For a solution of \eqref{four:eq:system},
\eqref{four:eq:boundary-identity} immediately gives
$\widehat H=-e^{-t}\sqrt d\,N_B<0$.

We specify the entire homotopy because its boundary modifications and
strictly increasing height term will be used in the a priori estimates.
The stages run from the desired system to the terminal problem with
solution $f=Z=0$; Section~\ref{four:sec:existence} transfers existence back
along this homotopy.  Take a smooth
$j_0:\R\to[0,1]$ which is zero for $t\le0$ and one for $t\ge1$.
For $0\le\lambda\le1$ put
\[
 V_\lambda=u\Achi\bigl(3\grad Z+\lambda K(w,\cdot)\bigr).
\]
In stages 1--3 use
\begin{equation}\label{four:eq:homotopy-flux}
\begin{cases}
 \divg V_\lambda=u\Xi &\text{in }\Omega_R,\\
 (V_\lambda)_\nu/u=
 [1-(1-\lambda)j_0(t)]
 [\cB-(1-\lambda)(\Achi K(w,\cdot))_\nu]
 &\text{on }B.
\end{cases}
\end{equation}
Keep the outer data $f=Z=0$ throughout.

Here are explicit smooth choices of the auxiliary trace terms.  Fix
$0<4\xi_1<b_+-b_-$ and smooth height switches $q_-,q_+$ such that
\[
\begin{array}{lll}
 q_-(z)=1 &(z\le b_-+\xi_1),&q_-(z)=0\quad(z\ge b_-+2\xi_1),\\
 q_+(z)=0 &(z\le b_+-2\xi_1),&q_+(z)=1\quad(z\ge b_+-\xi_1),
\end{array}
\]
where $q_-$ is nonincreasing and $q_+$ is nondecreasing.  Choose a smooth
direction switch $q_d$ with values in $[0,1]$, equal to one on
$(-\infty,-3/4]$ and zero on $[-1/2,\infty)$.  On each collar write
$\nu_s=\grad s/|\dd s|$, so $\nu_s=\nu$ at the face, and choose a smooth cutoff $\omega_B$ supported
there and equal to one near the face.  The collars are disjoint.  With
fixed constants
\[
 A_0>2\sup_B|H|+1,\qquad A_1>\sup_B|H|+1,
\]
define, extending by zero outside the collars,
\begin{equation}\label{four:sys:switches}
\begin{aligned}
 D_0(x',w,z)
  &=-A_0\omega_B(x')q_d(w\cdot\nu_s)q_-(z)
                      &&\text{on black collars},\\
 D_0(x',w,z)
  &= A_0\omega_B(x')q_d(-w\cdot\nu_s)q_+(z)
                      &&\text{on white collars},\\
 D_1(x',w)&=A_1\omega_B(x')\,w\cdot\nu_s
                      &&\text{on every collar}.
\end{aligned}
\end{equation}
Both terms vanish at $w=0$, are bounded for $|w|\le1$, and have bounded
derivatives of every order on bounded height ranges.  These bounds are
independent of $N$.  In addition $\partial_zD_0\ge0$.

The successive trace prescriptions and inner Dirichlet data are as
follows; the indicated parameter moves in the displayed direction.
\begin{enumerate}[label=\textup{Stage \arabic*:},leftmargin=*]
\item Decrease $\lambda$ from one to zero, keeping
      $F=h+C$ and $h|_B=b$.
\item Keep $\lambda=0$, increase $\alpha_0$ from zero to one, and use
      $F=h+C+\alpha_0D_0(x',w,h)$ and $h|_B=b$.
\item Keep $\lambda=0$, increase $\zeta$ from zero to one, and use
\begin{equation}\label{four:sys:trace-prescriptions}
\begin{aligned}
 F={}&h+(1-\zeta)\bigl[C+D_0(x',w,h+\zeta b(x'))\bigr]\\
    &+\zeta\bigl[\tr_{\Achi}K+D_1(x',w)\bigr],
 \qquad h|_B=(1-\zeta)b.
\end{aligned}
\end{equation}
\item Retain the terminal prescription
      $F=h+\tr_{\Achi}K+D_1$ and $h|_B=0$.  Keep $\lambda=0$ and
      multiply both right sides of \eqref{four:eq:homotopy-flux} by a common
      parameter $\gamma$, decreasing from one to zero.
\end{enumerate}
The endpoints agree, so this is a continuous homotopy.  All trace right
sides have derivative at least one with respect to $h$, with $(x',w,t)$
fixed.  This is also the relevant derivative when comparing scalar
solutions at a shared gradient and shared $Z$, since $t$ is independent
of the value of $f$.

Two consequences of these choices will be useful.  First, at a stage 3
face, evaluate the prescribed right side at $h=(1-\zeta)b$ and at the
limiting values $w=\pm\nu$, $d=\chi=0$.  It obeys
\begin{equation}\label{four:sys:direction-inequalities}
 F(\nu)\ge P_B+H,\qquad F(-\nu)\le P_B-H.
\end{equation}
Indeed, on a black face $b+C=P_B+H$, $D_0(\nu,b)=0$, and
$D_0(-\nu,b)=-A_0$; on a white face $b+C=P_B-H$,
$D_0(-\nu,b)=0$, and $D_0(\nu,b)=A_0$.  The terminal expressions are
$P_B\pm A_1$.  Each inequality follows for both endpoints from the
choices of $A_0,A_1$, and hence for their convex combination.

Second, throughout stage 4 the scalar solution is $f=0$, for every trial
$Z$.  Its equation is
$\tr_{\Achi}H^f=\eta_N f+D_1$, with zero boundary values.  At a
positive interior maximum $w=0$, $D_1=0$, $\Achi=\Id$, and
$l\Delta f=\eta_N f>0$, a contradiction; a negative minimum is
excluded in the same way.  Consequently in this stage
\begin{equation}\label{four:sys:terminal-system}
 f=0,\quad t=Z,\quad
 \cT=\cT_K+3(p+1)|\dd t|^2,\qquad
 \cT_K=\tfrac12(|K|^2+\tau^2),
\end{equation}
and its two remaining equations are exactly
\[
 \divg(3u\grad t)=\gamma u\Xi,\qquad
 3\partial_\nu t=\gamma[1-j_0(t)](-H-N_B).
\]

\subsection{Coarse heights and an outer boundary estimate}

\begin{lemma}[Height bound before the floor]\label{four:lem:coarse-height}
There exist fixed $r_0$ and $C_*>0$, independent of $N$, $R$, and all
homotopy parameters, with the following property.  Every smooth solution
of the trace equation in any stage, with its prescribed Dirichlet data,
on $\Omega_R$ for $R\ge2r_0$ satisfies
\begin{equation}\label{four:sys:coarse-height}
 |h|+|F|\le C_*,\qquad
 |h|\le C_*(r^{-3/2}-R^{-3/2})\quad(r_0\le r\le R).
\end{equation}
Here $Z$ may be arbitrary; the divergence equation and the condition
$t\ge-\eps$ are not assumed.  In particular, with the prescribed outer
value $Z=0$,
\begin{equation}\label{four:sys:outer-gradient}
 |\grad f|\le C_*\eta_N^{-1}R^{-5/2}\quad\text{on }S_R.
\end{equation}
For every fixed $N$ there is $R_{\min}(N)\ge2r_0$, chosen to tend to
infinity with $N$, such that every such solution for $R\ge R_{\min}(N)$
has
\begin{equation}\label{four:sys:outer-floor}
 t> -\eps\quad\text{on }S_R.
\end{equation}
\end{lemma}

\begin{proof}
At an interior extremum of $h$, $w=0$ and both auxiliary terms vanish.
The trace equation becomes
\[
 l\Delta f=h+C-\tau\quad\text{in stages 1 and 2},\qquad
 l\Delta f=h+(1-\zeta)(C-\tau)\quad\text{in stage 3}.
\]
Since $l>0$, comparison at a maximum and a minimum bounds $|h|$ by
$\max(\sup_B|b|,\|C-\tau\|_\infty)$.  Stage 4 has $h=0$.
All prescribed trace right sides then bound $|F|$ by a fixed constant,
because $D_0,D_1$ are bounded and $\|\Achi\|\le1$.

Choose $r_0$ beyond the supports of $K,C,b,D_0,D_1$.  At a comparison
with a radial function, the axis of $\Achi$ is
$\grad_g r/|\dd r|_g$.  Uniformly for $0<\chi\le1$, the prepared end
expansion gives
\begin{equation}\label{four:sys:radial-trace}
 \tr_{\Achi}\Hess_g r^{-3/2}
 =\frac32\left(-3+\frac52\chi\right)r^{-7/2}
                +O(r^{-11/2})<0
\end{equation}
after increasing $r_0$.  The leading coefficient is at most $-3/4$,
so the choice is independent of $N$.  Take $C_*$ large enough that
$C_*(r_0^{-3/2}-R^{-3/2})$ dominates the global height bound for every
$R\ge2r_0$.  On this annulus the scalar equation is
\[
 \frac{l}{\eta_N\sqrt D}\tr_{\Achi}\Hess_g h=h.
\]
The functions $\pm C_*(r^{-3/2}-R^{-3/2})$ are respectively upper and
lower barriers by \eqref{four:sys:radial-trace}.  At a putative first
comparison point the gradients coincide, so the same positive trace
matrix and the same positive prefactor occur in both expressions.
The Hessian order and the proper height sign give a contradiction.
This proves the end estimate.  Differentiating the comparison at the
outer boundary, where $h=0$ tangentially, gives
\eqref{four:sys:outer-gradient}, absorbing the uniform end normal comparison
into $C_*$.

Finally let $G(t)=t+\frac16\log(1+l(t)^2\sigma^2)$ at an outer point.
It is strictly increasing and $G(t)=Z=0$.  Thus $t>-\eps$ if
\[
 G(-\eps)=-\eps+\frac16\log(1+\ell^2\sigma^2)<0.
\]
By \eqref{four:sys:outer-gradient}, it suffices to choose
$R_{\min}(N)$ so large that
$C_*^2\ell^2\eta_N^{-2}R_{\min}(N)^{-5}<e^{6\eps}-1$.
This uses only fixed-$N$ enlargement of the outer radius.
\end{proof}

\subsection{The stationary comparison at a minimum of \texorpdfstring{$t$}{t}}

The scalar identity is useful at a floor contact because a second,
Riemannian Gauss calculation gives an upper bound for the same scalar
curvature.  Define, for a symmetric tensor $A$ and a chosen unit axis $e$,
\begin{equation}\label{four:sys:gauss-quadratic}
 \mathcal S(A)=\frac13(\tr_\perp A)^2
       -\frac12|A_\perp^0|^2
       +d A_{ee}\tr_\perp A-d|A_{e\perp}|^2.
\end{equation}

\begin{lemma}[Stationary Gauss comparison]\label{four:sys:gauss-lemma}
At an interior local minimum of $t$,
\begin{equation}\label{four:sys:gauss-bound}
 \frac12e^{2t}\Scal_{\ghat}
 \le\frac12\Scal_g-v\Ric_g(e,e)+\mathcal S(H^f).
\end{equation}
At every point where $\dd t=0$ there are an absolute $c_2>0$ and a
polynomial bound $\Pi_N$ such that
\begin{equation}\label{four:sys:floor-coercivity}
 \cT-\mathcal S(H^f)
          \ge c_2\cT-\Pi_N\bigl(vF^2+|K|^2\bigr).
\end{equation}
These conclusions extend to $\dd f=0$ with any unit axis.
\end{lemma}

\begin{proof}
Use the Riemannian warped product $g+l^2(\dd b')^2$ and its graph
$b'=f$.  At $\dd t=0$ one has $\dd l=0$, so the graph second form is
$H^f$, up to its immaterial overall sign.  Its inverse metric is $\Ad$.
The contracted second-form contribution is
\[
 \tfrac12\{(\tr_{\Ad}H^f)^2-|H^f|_{\Ad\otimes\Ad}^2\}
   =\mathcal S(H^f),
\]
as is seen by writing $\Ad=\operatorname{diag}(1,1,1,d)$.
The ambient scalar curvature and Ricci blocks are
\[
 \Scal_g-2\Delta l/l,\qquad
 \Ric_g-\Hess l/l,\qquad -\Delta l/l,
\]
for the scalar, horizontal block, and unit vertical direction;
the mixed block is zero.  The graph unit normal has horizontal part
$-w$ and vertical component $\sqrt d$.  Therefore the ambient scalar
minus twice its normal Ricci, divided by two, is
\[
 \frac12\Scal_g-v\Ric_g(e,e)
                     -v\tr_\perp(\Hess l/l).
\]
At a stationary point $\Hess l/l=p\Hess t$, and the conformal term
is $-3\tr_{\Ad}\Hess t$.  Thus we have, more precisely, the equality
\begin{equation}\label{four:sys:stationary-gauss-exact}
\begin{aligned}
 \frac12e^{2t}\Scal_{\ghat}
 ={}&\frac12\Scal_g-v\Ric_g(e,e)+\mathcal S(H^f)\\
 &-(3+pv)\tr_\perp\Hess_g t-3d\Hess_g t(e,e).
\end{aligned}
\end{equation}
At a minimum the last two terms are nonpositive, proving
\eqref{four:sys:gauss-bound}.

For \eqref{four:sys:floor-coercivity}, first subtract $\mathcal S(S_f)$.
An expansion using $\tr T=F-\chi j$ gives exactly
\begin{equation}\label{four:sys:stationary-difference}
\begin{aligned}
 \cT-\mathcal S(S_f)
 ={}&|T^0|^2+(1+d)|M|^2
       +\chi(1+d-2\chi/3)j^2\\
 &+(\chi/3-d)Fj+F^2/3.
\end{aligned}
\end{equation}
At $d=\chi=1$ the $(F,j)$ quadratic is
$F^2/3-2Fj/3+4j^2/3$, whose matrix has determinant $1/3$ and positive
trace.  It is uniformly positive in a fixed neighborhood of this point.
Since
\[
 1-d=v,\qquad 1-\chi=\frac{(3+p)v}{3+pv},
\]
there is an absolute $\kappa_0>0$ such that this neighborhood contains
all inputs with $(1+p)v\le\kappa_0$.  In that region
\eqref{four:sys:stationary-difference} controls an absolute fraction of
$|T^0|^2+|M|^2+\chi j^2+F^2$.

In the remaining region, use
$\chi(1+d-2\chi/3)\ge\chi$ and $|\chi/3-d|\le d$ to absorb the
mixed term into, for example, $\chi j^2/4$ and a multiple of
$(d^2/\chi)F^2$.  The only ratio required is
\begin{equation}\label{four:sys:polynomial-ratios}
 \frac{d^2}{\chi}=\frac{d(3+pv)}3\le\frac{3+N}{3},\qquad
 v>\frac{\kappa_0}{1+p}\ge\frac{\kappa_0}{1+N}.
\end{equation}
Keeping also a fixed fraction of $F^2$ at the cost of another multiple
of $F^2$, this proves a lower bound by an absolute fraction of the
stationary norm minus $C(1+N)^2vF^2$.  Equation
\eqref{four:sys:stationary-coercivity} converts the norm to $\cT$.

It remains to replace $S_f$ by $S_f-K$.  Polarizing
\eqref{four:sys:gauss-quadratic} bounds the difference by
\[
 C\bigl(|T|+d|j|+d|M|\bigr)|K|+C|K|^2.
\]
The stationary norm controls $|T|$ and $|M|$, while
$d|j|\le(d^2/\chi)^{1/2}(\chi j^2)^{1/2}$.
By \eqref{four:sys:polynomial-ratios} and Young's inequality, for every
fixed $\alpha>0$ this is at most
$\alpha\cT+C_\alpha(1+N)|K|^2$.  Choose $\alpha$ to retain half
the already obtained positive coefficient of $\cT$.  This proves
\eqref{four:sys:floor-coercivity} with absolute $c_2>0$.
\end{proof}

\begin{lemma}[Errors from the trace switches and drift]
\label{four:sys:stationary-errors}
Let $\mathcal K$ be a fixed compact enlargement of the supports of
$K,C,b,D_0,D_1$.  For every fixed $\delta_2>0$, every smooth trace
solution in stages 1--3 satisfies, at a point where $\dd t=0$,
\begin{equation}\label{four:sys:derivative-errors}
\begin{aligned}
 w(F)&\ge\eta_N a\sigma-\delta_2\cT-\Pi_N\mathbf1_{\mathcal K},\\
 u^{-1}\bigl|\divg(u\Achi K(w,\cdot))\bigr|
   &\le\delta_2\cT+\Pi_N\mathbf1_{\mathcal K}.
\end{aligned}
\end{equation}
The polynomial may depend on $\delta_2$ and the fixed geometric choices,
but no negative power of $\ell$ or $\eta_N$ is required.
\end{lemma}

\begin{proof}
At $\dd t=0$, differentiation gives
\[
 \nabla_iw=\Ad H^f_{i\cdot},\qquad
 \dd\log u=H^f(w,\cdot),\qquad \dd p=p'(t)\dd t=0.
\]
These formulas and \eqref{four:sys:stationary-coercivity} show
\begin{equation}\label{four:sys:stationary-derivatives}
 |\nabla w|+|\dd\log u|_{\Achi}+|\nabla\Achi|
                  \le\Pi_N(\sqrt{\cT}+|K|).
\end{equation}
For clarity, the axial entry in $\nabla w$ is $dH^f_{ee}$, and its
control uses $d^2/\chi\le(3+N)/3$.  In $|\dd\log u|_{\Achi}$
the axial entry has the factor $\sqrt\chi$.  Differentiating
$\Achi=\Id-(3+p)w\otimes w/(3+pv)$ uses only the preceding
$\nabla w$ and coefficients bounded by powers of $1+N$, since
$3+pv\ge3$.  This proves \eqref{four:sys:stationary-derivatives} without
dividing by $l$ or $\eta_N$.

Write the relevant prescribed trace as $\Phi(x',w,t,h)$.  Its height
derivative is at least one, so the chain rule contains the favorable term
\[
 \Phi_h w(h)\ge w(h)=\eta_N a\sigma.
\]
All other terms are supported in $\mathcal K$.  Lemma~\ref{four:lem:coarse-height}
bounds their height arguments in a fixed range; derivatives of $D_0,D_1$
and $b$ are consequently bounded.  Derivatives of $\tr_{\Achi}K$ use
$\nabla K$ and \eqref{four:sys:stationary-derivatives}.  The $t$ derivative
has zero contribution at the point in question.  Thus the remaining
absolute error is bounded by $\Pi_N(1+\sqrt{\cT})\mathbf1_{\mathcal K}$.
Young's inequality proves the first assertion.

Expanding the divergence in the second assertion differentiates
$\Achi$, $K$, and $w$, and adds
$\inner{K(w,\cdot)}{\dd\log u}_{\Achi}$.
The latter is bounded by $C|K|\,|\dd\log u|_{\Achi}$.
Equation~\eqref{four:sys:stationary-derivatives}, compact support, and Young's
inequality give the second assertion with the same allowed dependencies.
\end{proof}

\subsection{Exclusion of first contact with the floor}

\begin{proposition}[First-floor exclusion]\label{four:prop:floor-exclusion}
There are a fixed $\delta_0>0$ and a positive polynomial $C_N$ in
\eqref{four:sys:source} such that, for every $N\ge4$ and every
$R\ge R_{\min}(N)$, no smooth solution at any stage of the homotopy
satisfies
\[
 \min_{\overline\Omega_R}t=-\eps.
\]
The choices depend only on the prepared data, $\eps$, and the fixed
switches and collars.  They do not use a height bound sharper than
Lemma~\ref{four:lem:coarse-height}.
\end{proposition}

\begin{proof}
Suppose first that a floor minimum is interior and the parameter is in
one of stages 1--3.  At that point $\dd t=0$ and $\Hess t\ge0$.
Subtract \eqref{four:sys:gauss-bound} from \eqref{four:eq:scalar-identity} and
use $V_\lambda=V-(1-\lambda)u\Achi K(w,\cdot)$.  This gives
\begin{align*}
 u^{-1}\divg V_\lambda\ge{}&
 \mu-\tfrac12\Scal_g+J(w)+v\Ric_g(e,e)
        +\cT-\mathcal S(H^f)+w(F)-F\tau\\
 &-(1-\lambda)u^{-1}\divg(u\Achi K(w,\cdot)).
\end{align*}
Choose $\delta_2>0$ in Lemma~\ref{four:sys:stationary-errors} so small that
the two losses $\delta_2\cT$ leave at least $c_2\cT/2$ in
\eqref{four:sys:floor-coercivity}.  The remaining tensor expressions
$\mu-\Scal_g/2=(\tau^2-|K|^2)/2$, $J$, and $F\tau$ are compactly
supported and bounded by fixed constants, using
Lemma~\ref{four:lem:coarse-height}.  On the complement of $\mathcal K$,
$F=h$ and $F^2\le C r^{-3}$; also
$|\Ric_g|\le C r^{-4}$.  On $\mathcal K$ the positive weight
$\rho_0$ has a fixed positive minimum.  Hence
\[
 \mathbf1_{\mathcal K}+vF^2+v|\Ric_g|
                   \le C(\rho_0+v\rho_1).
\]
Combining these bounds yields an absolute $c_3>0$ and a polynomial
$Q_N\ge0$ such that, at the proposed minimum,
\begin{equation}\label{four:sys:floor-divergence-lower}
 u^{-1}\divg V_\lambda
   \ge c_3\cT+\eta_N a\sigma-Q_N(\rho_0+v\rho_1).
\end{equation}
Fix
$0<\delta_0<\min(c_3/2,1/4)$.
Since $t=-\eps$ gives $m_0(t)=1$, choose $C_N$ to exceed
$Q_N+1+\sup_\Omega(\rho/\rho_0)$.  The right side of
\eqref{four:sys:floor-divergence-lower} then strictly exceeds
$\delta_0\cT+\rho+\delta_0\eta_N a\sigma-C_N(\rho_0+v\rho_1)
=\Xi$, contradicting \eqref{four:eq:homotopy-flux}.  This choice is
polynomial in $N$.

At an interior minimum in stage 4, \eqref{four:sys:terminal-system} gives
$f=0$, $t=Z$, and $v=0$.  If $\gamma>0$, its equation at the
stationary point is
\[
 3\Delta t=\gamma\bigl(\delta_0\cT_K+\rho-C_N\rho_0\bigr).
\]
Enlarge $C_N$, still polynomially, so that
$C_N>1+\sup_\Omega(\delta_0\cT_K+\rho)/\rho_0$.
The right side is strictly negative, contradicting $\Delta t\ge0$.
This additional bound is finite because $K$ is compactly supported.

Consider next an inner boundary minimum in stages 1--3.
There $j_0(-\eps)=0$, so the same drift correction appears on both
sides of \eqref{four:eq:homotopy-flux}.  Cancelling it gives $V_\nu/u=\cB$.
Equation~\eqref{four:eq:boundary-identity} and $d>0$ imply
\[
 3\partial_\nu t=-H-N_B<0.
\]
This contradicts the nonnegative derivative into $\Omega_R$ at a
boundary minimum.  In stage 4 with $\gamma>0$, the corresponding exact
formula from \eqref{four:sys:terminal-system} is
$3\partial_\nu t=\gamma(-H-N_B)<0$, giving the same contradiction.
The outer boundary is excluded by \eqref{four:sys:outer-floor}.

Finally, at $\gamma=0$, $f=0$ and the remaining problem is
$\divg(3u\grad t)=0$ with homogeneous inner conormal data and $t=0$
on $S_R$.  Testing by $t$ gives $\int_{\Omega_R}3u|\dd t|^2=0$.
Connectedness and the nonempty Dirichlet boundary imply $t=0$.
Thus no stage admits floor equality.
\end{proof}

The order of choices is now fixed: choose the background geometric
data and $\eps$, then $\delta_0$ and the polynomial penalty $C_N$.
At each $N$, enlarge the outer radius as in
Lemma~\ref{four:lem:coarse-height}.  Subsequent estimates may increase the
lower bound on $N$ or $R_{\min}(N)$, but none will replace $C_N$ by an
arbitrary fixed-$N$ constant.  In particular, the first-floor argument
has paid for $F\tau$ using a bounded compact error and has not used the
eventual small-height absorption in \eqref{four:eq:trace-slack}.
